% TOP_PROBLEM: {"id": 20002543, "title": "Matrix virial decomposition and a collision-free tensor stress test", "category_id": 16, "category": {"id": 16, "name": "physics", "display_name": "Mathematical Physics", "description": "Problems at the intersection of mathematics and physics.", "slug": "physics", "order_index": 16, "created_at": "2026-07-31T15:26:25.671Z"}, "status": "partially_solved", "statusQualification": "The source has both a scalar conjecture and a proposed tensor variant. A tensor result is not a proof of the scalar conjecture, and dimensional/mass hypotheses must be stated for any claimed theorem. The catalogue title/status was paired with a machine-generated research attempt; that attempt is not peer-reviewed evidence of settlement. The statement covers the original source questions rather than only the title’s subcase."}
% FIRST_POSTED: 2026-10-11
% ENTRY_KIND: statement_only
% SUBMISSION: {"kind": "statement_only", "category": "top_problems", "problem_id": "20002543", "proposal_type": "established", "source": "American Institute of Mathematics, Variational Methods in Celestial Mechanics; notes by Kuo-Chang Chen, version 22 June 2003, source item 14. https://aimath.org/WWN/varcelest/varcelest.pdf", "status_evidence": "Inherited source-catalogue status partially_solved preserved. Background has machine-generated research claims explicitly marked not peer reviewed; these are not independent status evidence. Primary source accessed 11 October 2026; compared against the supplied catalogue excerpt. Historical wording and partial-result remarks were checked for scope; no comprehensive current-literature or proof-verification audit was performed.", "status_checked_on": "2026-10-11", "duplicate_check": "Full 15,773-record source registry checked for identical normalized statements; available site registry checked for identical normalized title/statement and canonical-ID collision. Canonical ID is absent from the site registry. Related canonical records 20002547. No mathematical-equivalence literature audit performed.", "attribution": "Alain Chenciner. Statement preparation: OpenAI GPT-6 Codex assistant, generation 11 October 2026; source-matching and notation editing by the assistant, no independent human mathematical verification."}
% SOURCE_STATEMENT: "Problem 14. Is there a conceptual proof for Saari's conjecture? Why not fix the moment of inertia tensor and ask the same question (maybe in higher dimensions)? (Alain Chenciner)"
% Preparation and final self-audit: CONTRIBUTING.md prompts C and D; source/audit artifacts retained in the local physics_statements_2026_10_11 batch.
% Imported from: physics_statements.tex; UnsolvedMath ID 20002543
% !TeX program = lualatex
\documentclass[11pt]{article}
\usepackage[margin=25mm]{geometry}
\usepackage{fontspec}
\setmainfont{Latin Modern Roman}
\usepackage{amsmath,amssymb,mathtools}
\usepackage{unicode-math}
\setmathfont{Latin Modern Math}
\usepackage{xurl,hyperref}
\hypersetup{colorlinks=true,urlcolor=blue}
\setcounter{secnumdepth}{0}
\setlength{\parindent}{0pt}
\setlength{\parskip}{5pt}
\setlength{\emergencystretch}{3em}
\newcommand{\sref}[2]{\hyperlink{source:#1:#2}{[S#2]}}
\newcommand{\cataloguescope}[1]{\paragraph{Recorded target.}#1}
\newcommand{\R}{\mathbb R}
\newcommand{\C}{\mathbb C}
\newcommand{\N}{\mathbb N}
\newcommand{\Z}{\mathbb Z}
\newcommand{\Q}{\mathbb Q}
\newcommand{\F}{\mathbb F}
\newcommand{\eps}{\varepsilon}
\begin{document}
% UnsolvedMath ID 20002543; source catalogue code AIM-PHYSICS-0051
\setcounter{section}{20002542}
\section{Matrix virial decomposition and a collision-free tensor stress test}\label{20002543}
{\small UnsolvedMath ID 20002543 · Mathematical Physics}
\subsection{Definitions and mathematical statement}

\paragraph{Shared notation.}$\mathbb N=\{1,2,\ldots\}$; the standard number fields have their usual meanings. All additional conventions are stated locally.


\paragraph{Definitions and notation.}In the Newtonian $n$-body problem, positions $q_i\in\mathbb R^d$ and masses $m_i>0$ obey $m_i\ddot q_i=\nabla_{q_i}U$, where $U(q)=\sum_{i<j}m_im_j/|q_i-q_j|$ and the gravitational constant is normalized to one. Here $T=\tfrac12\sum_i m_i|\dot q_i|^2$, $E=T-U$, and a collision has $q_i=q_j$. The center of mass can be fixed at zero; moment of inertia is $I=\sum_i m_i|q_i|^2$. A central configuration has $\nabla_{q_i}U=-\lambda m_iq_i$ in center-of-mass coordinates for a common scalar $\lambda$. After fixing $I=1$, it is a critical point of $U$ on the collision-free inertia ellipsoid. Congruent configurations and scales are factored out for a finiteness count. A relative equilibrium is a rigidly rotating solution with time-independent shape. Saari’s conjecture in its moment-of-inertia formulation asks whether constant scalar $I$ along a Newtonian solution forces a relative equilibrium. The inertia tensor in center-of-mass coordinates is the matrix $Q=\sum_i m_iq_iq_i^{\mathsf T}$ (or its usual trace-complement physical convention); fixing it is a distinct, stronger geometric condition. A tensor convention matters in higher dimension.

\paragraph{Catalogue wording retained for traceability.}
Give a conceptual proof of Saari’s conjecture. Also fix the moment-of-inertia tensor and investigate the corresponding question, potentially in higher dimensions.
\paragraph{Scope and source qualification.}The source has both a scalar conjecture and a proposed tensor variant. A tensor result is not a proof of the scalar conjecture, and dimensional/mass hypotheses must be stated for any claimed theorem. The catalogue title/status was paired with a machine-generated research attempt; that attempt is not peer-reviewed evidence of settlement. The statement covers the original source questions rather than only the title’s subcase.
\subsection{Short English statement}
Find a conceptual proof of Saari’s conjecture and investigate the fixed-inertia-tensor analogue.
\subsection{Sources}
\begin{itemize}
\item[\hypertarget{source:20002543:1}{S1}] ULAM.AI and the UnsolvedMath Contributors, UnsolvedMath, record 20002543 (AIM-PHYSICS-0051). Catalogue curation under CC BY 4.0; definitions and English exposition added locally.
\url{https://huggingface.co/datasets/ulamai/UnsolvedMath}.
\item[\hypertarget{source:20002543:2}{S2}] American Institute of Mathematics, Variational Methods in Celestial Mechanics; notes by Kuo-Chang Chen, version 22 June 2003, source item 14.
\url{https://aimath.org/WWN/varcelest/varcelest.pdf}.
\end{itemize}
\label{end:20002543}
\end{document}
